\section{Parametrisation by observables and the expected loss}\label{sec:calibration}

The closed-form model is parametrised by observables, not fitted to price moments. Each primitive maps to a market quantity: the effective primary-redemption access $R$ to reserve liquidity and the issuer's primary arrangements, the price impact $\psi$ to secondary market depth, the signal dispersion $\sigma$ to the breadth of the holder base, and the fundamental volatility $\sigma_\theta$ to the frequency of stress. The counterparty parameters $(h,\xi)$ are read from market proxies as in Section~\ref{sub:counterparty}. Table~\ref{tab:cg_params} reports the values for the two coins.

\begin{table}[H]
\centering
\caption{Observable parametrisation of the closed-form model. $\psi$ is backed out from the observed episode trough $1-\psi(1-R)$; $R$ is the effective primary-redemption access in stress; $\sigma_\theta$ is set so the endogenous stress probability matches the empirical per-coin frequency.}\label{tab:cg_params}
\begin{tabular}{llcccccc}
\toprule
Coin & Episode & $R$ & $\psi$ & $\sigma$ & $\sigma_\theta$ & $h$ & $\xi$ \\
\midrule
USDC & SVB, Mar 2023   & 0.50 & 0.196 & 0.22 & 0.0074 & 0.01 & 0.20 \\
USDT & Terra, May 2022 & 0.85 & 0.166 & 0.24 & 0.0507 & 0.02 & 0.30 \\
\bottomrule
\end{tabular}
\end{table}

The single structural difference that drives the two coins apart is $R$, the access to par redemption when the shock hits. USDC's primary redemption was constrained over the March 2023 banking weekend, when the reserve held at Silicon Valley Bank was frozen, so its effective $R$ was low and the secondary absorbed the conversions, taking the price to 0.90. USDT's primary redemption stayed open through the May 2022 Terra fallout, about USD 8 billion clearing at par, so its effective $R$ was high and the price dipped only to 0.975. The backed-out price impacts are similar (0.20 and 0.17), so the depth of the two secondary markets is comparable; the difference in realised depeg is a difference in primary-redemption access, exactly the channel Proposition~\ref{prop:reserves} prices.

\subsection{The endogenous stress probability matches the data}\label{sub:ps_validation}

The stress probability $p_s$ of Proposition~\ref{prop:ps} is endogenous, a function of the threshold and the fundamental volatility. Evaluated at the parametrisation, it is $0.0052$ for USDC and $0.00031$ for USDT. These track the empirical per-coin frequencies of below-par stress hours over 2022--2026, roughly $0.005$ for USDC and $0.0003$ for USDT, computed directly from the hourly series. The model thus reproduces how often each coin is in stress from its reserve and depth primitives, rather than imposing a single frequency shared across coins. The reserve comparative static distinguishes the two: $\partial \Pi_{\text{risk}}/\partial R$ is $-0.56$ for the reserve-sensitive USDC and $-0.19$ for the counterparty-dominant USDT, so a liquidity requirement bites where the depeg risk is, and barely moves the coin whose risk is on the issuer's balance sheet.

\subsection{The model against a reduced-form benchmark}\label{sub:interp}

Because the counterparty and conversion pieces are observable, a natural question is what the coordination model adds beyond a purely reduced-form account. We answer it directly, on the full holding-and-conversion cost so the comparison has content. A reduced-form account built only from observables, the counterparty term $h\xi$, the off-ramp basis, and the empirical stress frequency times the observed episode depth, gives 43.6 basis points for USDC and 72.3 for USDT. The closed-form model gives 42 and 72 for the same object, the holding risk of Table~\ref{tab:dual_asset_premium} plus the conversion basis. The two agree, which is the point: the model does not inflate the number, and the cost is disciplined by observables. What the model adds is not the level but the mechanism, the endogenous per-coin stress probability of Proposition~\ref{prop:ps}, the reserve comparative static of Proposition~\ref{prop:reserves}, and the threshold that the reduced-form account cannot deliver and that the policy questions of Section~\ref{sec:discussion} require. The model's stress depeg is small in expectation because the deep trough is a tail event; the cost is dominated by the observable counterparty and conversion terms, and the model earns its place by pricing how those terms move with reserves, depth, and the regulation that governs them.
